\section{Design}
\label{sec:design}

\subsection{System Overview}

\namei treats a dynamic filesystem view as late binding of a pathname or
directory entry to an existing VFS object. Existing systems commonly couple
that binding either to namespace materialization before lookup or to a
filesystem implementation that owns lookup and subsequent operations. \namei
separates the binding policy from filesystem ownership: policy chooses a
bounded binding, while the VFS validates the result, owns the selected object,
and executes ordinary lower-filesystem semantics.

Three invariants organize the design. First, every VFS path that observes a
name implements the same bounded binding contract. Second, policy returns only
names or opaque handles; the kernel alone validates object type, scope, and
lifetime before installing a result. Third, a selected binding changes the
object on which ordinary VFS processing operates, not the semantics of those
operations. Permissions, file operations, data, writes, page cache,
persistence, and consistency remain with the VFS and lower filesystem. A
separate deployment requirement keeps this optional boundary off disabled and
unmanaged paths before policy-context construction.

The central design question is the acceptance contract for a policy-selected
object. Native name resolution does not return a bare inode: the surrounding
\code{nameidata} records how the object was reached, which operation is being
performed, which root and mount restrictions apply, whether the walk is under
RCU, and which sequence values protect it from concurrent namespace changes.
\namei may replace the visible object only when it can preserve the relevant
state or fail the operation before the result is used. This contract, rather
than the BPF invocation itself, determines the mechanism boundary.

Figure~\ref{fig:overview} shows the resulting flow. An application in a
workload scope performs an ordinary lookup, open, stat, exec, or readdir. The
VFS reaches the \namei decision point before committing to the visible object
for the current component or directory entry. The eBPF policy returns a bounded
path-view action, and ordinary VFS and lower-filesystem logic resumes after
kernel validation.

\begin{figure}[t]
\centering
\scriptsize
\fbox{\parbox{0.94\linewidth}{
\centering
\textbf{Linux VFS name-resolution machinery}\\[2pt]
\begin{tabular}{@{}c|c|c@{}}
component walk & final open/create & lower directory iterator \\
\end{tabular}\\[3pt]
static key and exact-parent filter: disabled or unrelated paths continue
without context construction\\[3pt]
$\downarrow$ managed parent\\[-2pt]
\fbox{\parbox{0.82\linewidth}{
\centering
one eBPF decision contract supplies policy for \code{LOOKUP} or
\code{READDIR}\\
\code{PASS}, \code{HIDE}, bounded component \code{REDIRECT}, or registered
target ID
}}\\[3pt]
$\downarrow$ kernel validation and object handoff\\[-2pt]
\begin{tabular}{@{}c|c|c@{}}
\begin{tabular}{c}\code{PASS}/\code{REDIRECT}\\stay in or convert\\the current walk\end{tabular} &
\begin{tabular}{c}\code{SELECT_TARGET}\\borrow registered path\\under RCU, then validate\end{tabular} &
\begin{tabular}{c}\code{READDIR}\\filter lower entries\\preserve iterator position\end{tabular} \\
\end{tabular}\\[3pt]
$\downarrow$\\[-2pt]
ordinary VFS permission, path-walk, open, and iterator completion\\[2pt]
$\downarrow$\\[-2pt]
lower filesystem retains dentries, inodes, data, writes, page cache,
persistence, and consistency
}}
\caption{\namei separates late-binding policy from filesystem ownership. BPF
chooses a bounded binding; the kernel validates and installs it according to
the surrounding component, final-open, readdir, and RCU/ref-walk path; ordinary
VFS semantics then operate on the selected object.}
\label{fig:overview}
\end{figure}

\subsection{One Binding Contract Across Resolution Paths}

Linux does not resolve every name through one function. Intermediate
components pass through the component walk, the trailing component of
\code{open()} has separate create and open handling, and \code{readdir()}
obtains names from a filesystem iterator. \namei applies one policy contract at
all three points rather than assigning each call site an independent rule
language. The input identifies a lookup or directory-entry event; policy state
and the current parent/name pair determine the result.

The action meanings are consistent across those paths. \code{PASS} exposes the
ordinary lower result. \code{HIDE} returns absence from lookup and suppresses
the corresponding directory entry. \code{REDIRECT} substitutes a validated
component name for lookup or for the emitted directory entry.
\code{SELECT_TARGET} selects a registered existing object during lookup;
directory iteration over a selected directory proceeds after that directory
has been opened. The current design does not use \code{SELECT_TARGET} to invent
directory entries, because synthetic aliases require an object model and
stable iterator positions beyond existing-entry filtering.

The ABI gives lookup and readdir one action vocabulary, but the kernel invokes
the two event types independently. A policy is responsible for making its
lookup result and directory view coherent; workload oracles test that
obligation. Target selection changes the object on which subsequent ordinary
VFS processing operates. Create semantics remain outside selection:
\code{REDIRECT} and \code{SELECT_TARGET} reject \code{LOOKUP_CREATE}, and the
policy does not decide target creation or copy-up behavior.

The contract linearizes each decision at one program invocation. It does not
provide a transaction across independent system calls or across all entries of
a directory stream. A controller can publish immutable target generations and
switch one policy-state selector atomically, as the evaluated workloads do,
but policies that require cross-name transactional updates need a broader
mechanism. Stating this update model separates binding coherence for a fixed
policy state from filesystem transaction semantics that \namei does not own.

\subsection{Accepting A Policy-Selected Object}

A decision is not yet a usable VFS object. BPF cannot return a kernel pointer
or an arbitrary path string, and cache-hot RCU-walk does not hold ordinary
references that a policy can transfer. \namei therefore separates policy
output from object ownership. A control plane registers an existing object,
the kernel retains its \code{struct path}, and policy selects it by an opaque
ID. The kernel validates the action, target type, scope, and operation before
installing the result in the current namei state.

Object handoff differs between ref-walk and RCU-walk without changing policy meaning.
Ref-walk acquires an independent path reference directly. RCU-walk borrows the
registry's kernel-held path, then uses ordinary namei sequence validation to
obtain stable references or restart if concurrent replacement invalidates the
walk. Policy chooses only the binding; object lifetime and the transition from
a speculative walk to stable VFS state remain kernel responsibilities.

The accepted subset is defined by five checks. First, registration from an
open descriptor establishes both controller authority and kernel-held target
lifetime; policy may name only the resulting cgroup-scoped ID. Second, a target
must be type-compatible with its position: an intermediate target is a
directory, while a final target is an existing regular file or directory.
Third, selection cannot own creation, so \code{LOOKUP_CREATE} is rejected.
Fourth, selection fails when an \code{openat2()} scoped walk could lose its
root constraint, or when \code{RESOLVE_NO_XDEV} would cross a mount. Fifth,
RCU-walk may borrow a registered path only until ordinary namei validation can
stabilize it; otherwise the walk restarts or fails. Unsupported cases do not
fall back to the unmodified lower name.

Registration is therefore an authorization boundary, not merely a lookup
cache. The controller authorizes an already opened object for selection by one
policy scope. Subsequent VFS processing checks the logical walk and selected
object, but does not replay traversal through the target's physical parents.
Policies that require unprivileged arbitrary-path selection or scoped traversal
through a replacement path need a broader mechanism.

\subsection{Preserving Filesystem Semantics}

After installing a selected object, \namei resumes the surrounding VFS state machine rather than
starting a second path walk or forwarding filesystem methods. Intermediate
directory selection continues component walking from the selected directory;
final selection continues permission, stat, access, open, or exec completion
on the selected object; opening a selected directory uses its lower iterator.
An already opened file or directory subsequently executes lower-filesystem
operations without reinvoking policy.

This preservation invariant is what keeps \namei from becoming a filesystem.
The extension does not emulate permission checks, mount handling, dentries,
inodes, reads, writes, page cache, persistence, or filesystem-specific
behavior. For any operation after a successful binding, the required result is
the result of ordinary VFS execution on the selected lower object.

\subsection{Name-Resolution Boundary Goals}

Four evaluation goals follow from this ownership split: expressiveness,
semantic preservation, placement cost, and fail-closed validation. Like
\code{sched_ext}, \namei lets BPF choose policy while the kernel keeps the VFS
mechanisms that own objects and execute filesystem semantics.

The design defines a narrow kernel extension for name-resolution policy while
leaving full filesystem expressiveness to FUSE, stackable, or custom
filesystems. Table~\ref{tab:req-design} maps the four design goals to the RQs.

\begin{table}[t]
\centering
\scriptsize
\begin{tabular}{@{}p{0.19\linewidth}p{0.5\linewidth}p{0.17\linewidth}@{}}
\hline
Goal & Design choice & RQ \\
\hline
G1 Expressiveness & invoke one workload-scoped decision contract at lookup and directory-enumeration time & RQ1 \\
G2 Semantic preservation & select only existing lower objects and resume ordinary VFS processing for data, writes, permissions, page cache, and persistence & RQ1, RQ3 \\
G3 Placement cost & run bounded policy at lookup/readdir instead of routing the same decision through a filesystem service & RQ2 \\
G4 Fail-closed validation & use verified eBPF outputs plus kernel validation so malformed or unsupported bindings fail visibly & RQ3 \\
\hline
\end{tabular}
\caption{Design-goal mapping. If an oracle needs synthetic contents, custom
metadata, or data-path mediation, the effect belongs outside \namei's
boundary.}
\label{tab:req-design}
\end{table}

\subsection{Fast-Path Contract}

The extension sits on a system-wide hot path, so its cost model is part of the
interface rather than an implementation afterthought. When no program is
attached, lookup and directory iteration must follow the ordinary VFS path
after one static-key test. When a program is attached to a limited set of
parents, an unrelated parent must be rejected before policy-context
construction or workload-scope discovery. This first test is conservative:
false positives may execute the authoritative scope check, but a false
negative would bypass policy and is forbidden.

For managed parents, cache-hot lookup may already be in RCU-walk. The decision
contract therefore cannot require a sleeping lock or an owned dentry/mount
reference. \namei separates the BPF decision from action-specific object
handling: pass-through can remain in RCU-walk, while an action that needs
refcounted state must use normal namei conversion and validation. Walk mode is
not exposed as workload policy, so the same pathname and policy state do not
produce different decisions merely because the VFS cache changed.

A successful in-place conversion applies the saved decision without calling
the program again. If conversion fails and ordinary namei returns
\code{ECHILD}, however, the VFS may restart the complete pathname walk and
execute policy again. The interface therefore provides at-least-once execution
across full VFS restarts, not exactly-once execution. Policies must keep
externally visible side effects idempotent; map reads and writes that define the
decision remain subject to the same restart rule.

\subsection{Policy Model}

The workloads require policy to see lookup and directory-enumeration events,
carry workload state, return only bounded path-view actions, and let the kernel
validate the result.

The policy input contains operation context, the current component, parent
context, and policy state. The policy state is workload-defined but not a
filesystem object model. For the representative workloads, policy state is a
branch or checkpoint identifier in an agent workspace, a cache or validation
epoch in an environment run, or a configuration/secret epoch in a service.

The output passes the lower lookup, redirects selection to a replacement
component, hides a name from lookup and directory enumeration, or selects a
registered lower object. Access-control denial remains outside the model
because \namei focuses on selecting or hiding existing objects during name
resolution. Directory enumeration uses the same decision function and action
vocabulary as lookup; the policy must implement the corresponding direction,
and evaluation checks the resulting lookup/readdir view.

The policy cannot return a kernel pointer or an arbitrary pathname. A control
plane first registers an open existing object, and the kernel retains its
\code{struct path} under an opaque ID. The policy may return only that ID.
Intermediate selections must identify directories; final selections may
identify existing regular files or directories. The selected object then
continues through normal namei validation, permission checks, open completion,
and lower-filesystem file operations.

The kernel also validates action type, name bounds, target identity, selected
object type, create restrictions, and path-walk scope. Invalid decisions fail
visibly rather than falling back to a weaker policy. This registered-object
model bounds what policy can select without turning the hook into arbitrary
recursive path rewriting.

\subsection{Filesystem Ownership Boundary}

\namei excludes filesystem ownership by leaving inode operations, file
operations, dentry and inode allocation, synthetic contents, custom metadata,
read/write mediation, distributed indexes, and cross-path transactions outside
the hook.

The name-resolution boundary helps only when the source transition can
be expressed as lookup selection or directory visibility. When the oracle
requires data-path mediation, synthetic content, custom metadata, or
application-level orchestration, the transition remains motivation or related
work rather than \namei evidence.
